\setcounterpageref{ex}{0}
\TaoPhieuKiemTra{05}

\setcounter{ex}{0}
\setcounter{bt}{0}

\textbf{\large I. PHẦN TRẮC NGHIỆM (3,0 điểm)}

{\itshape Khoanh tròn vào chữ cái đặt trước câu trả lời đúng. Mỗi câu đúng được 0,5 điểm.}

\begin{ex} % Câu 1
Số tròn chục liền sau số $50$ là
\choice
{$51$}
{\True $60$}
{$40$}
{$70$}
\loigiai{
Số tròn chục liền sau số $50$ là $60$.\\
\textbf{Đáp án đúng: B}
}
\end{ex}

\begin{ex} % Câu 2
Trong phép trừ có số bị trừ là $50$, số trừ là $20$. Hiệu là
\choice
{$70$}
{$20$}
{$40$}
{\True $30$}
\loigiai{
Hiệu là: $50 - 20 = 30$.\\
\textbf{Đáp án đúng: D}
}
\end{ex}

\begin{ex} % Câu 3
Quả dưa hấu cân nặng $4\text{ kg}$, quả bưởi cân nặng $2\text{ kg}$. Cả hai quả nặng
\choice
{$5\text{ kg}$}
{$7\text{ kg}$}
{\True $6\text{ kg}$}
{$2\text{ kg}$}
\loigiai{
$4 + 2 = 6\text{ kg}$.\\
\textbf{Đáp án đúng: C}
}
\end{ex}

\begin{ex} % Câu 4
Đoạn thẳng $MN$ dài $7\text{ cm}$, đoạn thẳng $PQ$ dài $10\text{ cm}$. Đoạn thẳng $PQ$ dài hơn đoạn $MN$ là
\choice
{\True $3\text{ cm}$}
{$2\text{ cm}$}
{$4\text{ cm}$}
{$17\text{ cm}$}
\loigiai{
$10 - 7 = 3\text{ cm}$.\\
\textbf{Đáp án đúng: A}
}
\end{ex}

\begin{ex} % Câu 5
So sánh: $45 + 23 \dots 70$.
\choice
{$>$}
{$=$}
{Không so sánh được}
{\True $<$}
\loigiai{
$45 + 23 = 68$. Vì $68 < 70$ nên điền dấu $<$.\\
\textbf{Đáp án đúng: D}
}
\end{ex}

\begin{ex} % Câu 6
Tính: $7 + 7 - 9 = \dots$
\choice
{$4$}
{\True $5$}
{$6$}
{$7$}
\loigiai{
$7 + 7 = 14$; $14 - 9 = 5$.\\
\textbf{Đáp án đúng: B}
}
\end{ex}

\vspace{4mm}
\textbf{\large II. PHẦN TỰ LUẬN (7,0 điểm)}

{\itshape Thí sinh trình bày chi tiết bài làm vào giấy kiểm tra.}

\begin{bt}[1,0 điểm]
Tính nhẩm:
\begin{multicols}{2}
\begin{enumerate}[a)]
  \item $9 + 4$
  \item $7 + 5$
  \item $12 - 6$
  \item $14 - 7$
\end{enumerate}
\end{multicols}
\loigiai{
\begin{itemize}
  \item $9 + 4 = 13$
  \item $7 + 5 = 12$
  \item $12 - 6 = 6$
  \item $14 - 7 = 7$
\end{itemize}
}
\end{bt}

\begin{bt}[1,5 điểm]
Đặt tính rồi tính:
\begin{multicols}{3}
\begin{enumerate}[a)]
  \item $33 + 47$
  \item $85 - 36$
  \item $51 + 46$
\end{enumerate}
\end{multicols}
\loigiai{
\begin{itemize}
  \item $33 + 47 = 80$
  \item $85 - 36 = 49$
  \item $51 + 46 = 97$
\end{itemize}
}
\end{bt}

\begin{bt}[1,0 điểm]
Tính giá trị biểu thức kèm đơn vị đo:
\begin{enumerate}[a)]
  \item $21\text{ kg} + 8\text{ kg} - 11\text{ kg}$
  \item $30\text{ }l - 11\text{ }l + 7\text{ }l$
\end{enumerate}
\loigiai{
\begin{enumerate}[a)]
  \item $21\text{ kg} + 8\text{ kg} - 11\text{ kg} = 29\text{ kg} - 11\text{ kg} = 18\text{ kg}$.
  \item $30\text{ }l - 11\text{ }l + 7\text{ }l = 19\text{ }l + 7\text{ }l = 26\text{ }l$.
\end{enumerate}
}
\end{bt}

\begin{bt}[1,5 điểm]
Bạn Giang có $20$ con tem. Bạn Hà có nhiều hơn bạn Giang $9$ con tem. Hỏi bạn Hà có bao nhiêu con tem?
\loigiai{
Số con tem bạn Hà có là:
\[ 20 + 9 = 29 \text{ (con tem)} \]
\textbf{Đáp số:} $29$ con tem.
}
\end{bt}

\begin{bt}[1,0 điểm]
Một thùng chứa đầy $40\text{ }l$ nước mắm. Buổi sáng cửa hàng đã bán ra $17\text{ }l$. Hỏi trong thùng còn lại bao nhiêu lít nước mắm?
\loigiai{
Trong thùng còn lại số lít nước mắm là:
\[ 40 - 17 = 23\text{ }l \]
\textbf{Đáp số:} $23\text{ }l$ nước mắm.
}
\end{bt}

\begin{bt}[1,0 điểm]
Tìm số có hai chữ số giống nhau mà tổng hai chữ số của nó bằng $14$.
\loigiai{
Vì hai chữ số giống nhau và có tổng bằng $14$ nên mỗi chữ số là $7$ (do $7 + 7 = 14$).\\
Vậy số đó là $77$.\\
\textbf{Đáp số:} $77$.
}
\end{bt}
